\section{A concrete path to complete Lean formalization}
\label{sec:lean}

\subsection{Repository interface and what has actually been checked}
The inspected ProveIt revision is
\begin{center}
\texttt{f7e5eac0a94a0e690a22cbeecc5f6cca9270e6df}.
\end{center}
At that revision the root toolchain file specifies Lean \texttt{v4.32.0},
and the root Lake configuration requests Mathlib \texttt{v4.32.0}.
The relevant source files are \texttt{ConciseIPRA/PR.lean},
\texttt{Deduction.lean}, \texttt{BiInterpretation.lean}, and
\texttt{NativeBiInterpretation.lean}~\cite{ProveIt}. We read their actual
interfaces. In particular the repository has the types
\[
 \texttt{PR k},\quad\texttt{PTm n},\quad
 \texttt{Fm PTm n},\quad\texttt{Prf PRAAx \,$\Gamma$\, $\varphi$},
\]
and the quantifier-free induction constructor
\texttt{PRAAx.ind} and derived rule \texttt{PRA.qfInd}.
Its theorem \texttt{pra\_iff\_native} translates an already proved
arithmetic sequent into \texttt{NPrf} with the translated formulas.

\textbf{Formalization status.} The human proofs and their explicit
finite estimates are the completed mathematical component. The supplied
Lean files are an \emph{uncompiled outline with definitions and small
candidate finite proof scripts}. There was no Lean, Lake, or Elan
executable in the preparation environment. No analytic theorem in this
package is claimed to be kernel checked, and no completed object-theory
proof certificate is supplied. Absence of \texttt{sorry} or added
\texttt{axiom} declarations in an outline is not evidence of completion.
The exact rational Python checks are separate implementation evidence.

\subsection{Definitions to formalize, in dependency order}
The file \texttt{FiniteDefinitions.lean} gives rational complex records,
finite sums and products, the regular-name and equality predicates,
apartness, the Bernoulli coefficient recursion, exponential and logarithm
polynomials, and the regularized Euler--Maclaurin expression. Its
rational-input branch at $s=1$ is decidable. For arbitrary fixed names,
replace that branch by the everywhere-defined integral or power series
$Q_N$ in \eqref{em:Q}; do not test exact equality of names.

A complete implementation should use the following modules. These names
are a proposed module map, not claims that such modules already exist.
\begin{longtable}{p{0.20\textwidth}p{0.72\textwidth}}
\toprule
Module & Required definitions and proved contract\\\midrule
\endfirsthead
\toprule Module & Required definitions and proved contract\\\midrule
\endhead
\texttt{Codes} & Signed integers, positive-denominator fractions, pairs,
finite lists, polynomial coefficient lists, bounded Boolean conjunction,
and decoding lemmas. All arithmetic must have explicit PR syntax.\\
\texttt{Names} & The predicates (F1)--(F2), arithmetic on names, rational
ball containment, supplied apartness, refinement and rational-gap lemmas.
Prove sum, product and inverse error rules.\\
\texttt{Elementary} & Fixed log/exp/sin/cos approximators using
\eqref{eq:exp-tail}--\eqref{eq:log-tail}; $\pi$ by arctangent; positive
square roots by bounded bisection. Prove their addition, derivative and
modulus identities from finite polynomials.\\
\texttt{FiniteCalculus} & The mesh sums and bound \eqref{eq:quad};
telescoping FTC, product rule, integration by parts, and substitutions on
certified compact intervals. Derivative moduli are data or explicit
consequences of the elementary formulas.\\
\texttt{Dirichlet} & The fixed right-half-plane evaluator; dyadic tail
\eqref{eq:d-tail}, derivative tail \eqref{eq:d-derivative-tail},
convolution \eqref{eq:d-convolution}, factorization, and the Euler,
M\"obius, eta and logarithm identities.\\
\texttt{BernoulliEM} & Finite coefficient recursions; endpoint,
reflection and difference identities; finite formula \eqref{em:finite}
with its negative integral remainder.\\
\texttt{Continuation} & The fixed evaluator $E_{m,N(k)}$, its $K_0/N$
error, order compatibility, regular part $H$, Taylor coefficients,
Stieltjes names, and the exact finite nonpositive values.\\
\texttt{GammaGateway} & The fixed reciprocal-gamma product, explicit
lower bounds, Euler integral, finite sine product, reflection and
normalized duplication, with their stated remainders.\\
\texttt{Theta} & Finite Fej\'er polynomials, the Gaussian transform,
periodization, the entire integral $A(s)$ and the pole-cleared expression
\textnormal{(TH8)}.\\
\texttt{Propagation} & Finite circle coefficient extraction,
\eqref{eq:analytic-propagation}, the backwards precision recurrence,
and its finite tree of sample points. This replaces an abstract
identity theorem.\\
\texttt{Hasse} & Finite differences $D_n$, the shifted finite Mellin
identity, the derivative-basis coefficient invariant, the bound
$C(n+1)^m$, and the outer geometric tail. Prove entire eta with finite
circle coefficient extraction and identify its cleared product with
zeta.\\
\texttt{Consequences} & Functional equations, $\xi$, Fourier Bernoulli
values, Mellin and cube integrals, the explicit boundary lower bound,
and the finite-block non-Cauchy witnesses.\\
\texttt{ObjectProofs} & For every compiled fixed evaluator, actual
\texttt{Prf PRAAx} derivations of its rational residual inequalities.
Only then use the existing native bridge.\\
\bottomrule
\end{longtable}

\subsection{The exact object-theory target}
A theorem about Lean's completed real or complex numbers, even if
constructively proved, is not itself an iPRA derivation. The target is a
finite rational predicate with the inputs coded as naturals. For example,
for a rational argument in an Euler--Maclaurin chart, a checker receives
\[
 u=\operatorname{code}(a,b,c,d,A,R,k,N,M,\ldots)
\]
where $(a/b)+i(c/d)$ is the input. It verifies the domain inequalities
and the rational finite-error conclusion. A failed domain test makes the
implication true; a valid input must satisfy the stated conclusion.
Return $0$ for success. Once this \emph{fixed} checker has been compiled
to a symbol $\mathsf{check}:\texttt{PR 1}$, the closed sentence is
\[
 \forall u\;\mathsf{check}(u)=0.
 \tag{L1}\label{eq:object-target}
\]
The required endpoint is a term of
\[
 \texttt{Prf PRAAx [] (residualSentence check)}.
 \tag{L2}\label{eq:object-proof}
\]
For a fixed nonrational input name, inline its fixed PR definition and
retain its regularity and domain formulas in $\Gamma$. This gives a
sequent rather than silently asserting that every name is regular.
Some conclusions, such as the non-Cauchy theorem, have an explicit
witness function: insert the given $M(N),L(N)$, prove the finite bound,
and then use intuitionistic existential introduction. There is no
unbounded minimization.

The included \texttt{ArithmeticTargets.lean} writes the syntax of
(L1)--(L2) and the final bridge. Its theorem
\texttt{native\_of\_completed} takes a completed object proof as its
hypothesis and transports it. It is an interface example, not a proof of
that missing hypothesis. The substantive unfinished formal work is to
construct the checkers, prove their coding meanings, and build their
\texttt{Prf PRAAx} terms from the finite arguments in the article.

\subsection{Quantifier-free invariants and precision accounting}
Here are the principal invariants, stated at the level at which they
should be compiled. Fix all exterior precision variables first.
\begin{enumerate}[leftmargin=*,itemsep=4pt]
\item For a finite sum: the accumulator equals the sum of the recorded
terms and its error is at most the sum of their assigned rational
radii. For a product, use the factor-replacement telescoping identity.
\item For a Bernoulli recursion: the stored coefficient vector has the
specified derivative vector, rational mean zero, and coefficient bound
$C_n$. Each identity is a bounded conjunction of rational equalities.
\item For finite Euler--Maclaurin: after $j$ cells, the accumulator is
exactly the endpoint correction plus the finite integral expression,
within the sum of the explicitly assigned quadrature errors. For the
integration-by-parts order, store the finite list of derivative and
Bernoulli coefficients and the remaining signed integral.
\item For a Dirichlet or Euler-product truncation: every included
factorization is unique and every omitted index lies in the declared
nonnegative tail majorant. This is finite arithmetic on bounded lists.
\item For gamma and theta: the accumulated rational evaluations lie
in their balls, with tail at most the explicit product, geometric or
exponential bound. No ``the series converges'' invariant is used.
\item For analytic propagation: build the finite tree of circle sample
points required by the successive coefficient quadratures. Store at each
node its rational residual bound and the backwards precision supplied
by (TH7). Tree depth and branching are PR numbers fixed before the
induction. The invariant quantifies only over this bounded coded tree.
\item For a final named equality: prove the finite residual inequality
for the already fixed precision, then introduce that precision
universally. Do not apply quantifier-free induction to equality of names.
\end{enumerate}
The cutoffs in the text are deliberately loose. Their role is to make
all existence witnesses primitive recursive and all induction invariants
finite; they are not assertions of efficient numerical complexity.
The Python evaluator uses tighter exact arithmetic whenever convenient,
but its reported enclosure always uses the proved analytic remainder.

\subsection{Executable evidence and reproducibility}
The program \texttt{code/zeta\_certificates.py} implements exact complex
rational balls with the $1$-norm, range-reduced logarithms, exponential
polynomials with rounded-ball repeated squaring, Bernoulli coefficients,
and the finite Euler--Maclaurin formula. Its remainder is
\[
 2\frac{C_{2m}\prod_{j<2m}\|s+j\|_1}{(2m)!}
       \frac{N^{1-\lfloor\Re s\rfloor-2m}}{\Re s+2m-1}.
 \tag{L3}\label{eq:code-remainder}
\]
The factor $2$ converts a Euclidean bound to a $1$-norm bound.
The exponent is made integral by rounding $\Re s$ down, which enlarges
the bound for $N\ge1$. When a Pochhammer factor is exactly zero at a
nonpositive integer, the remainder is exactly zero and rational integer
powers give the exact rational value.

The saved verification run passed $780$ exact polynomial
Euler--Maclaurin identities, including $300$ with nonzero remainder;
$359$ Bernoulli polynomial identities; $8$ reference Bernoulli values;
$13$ exact nonpositive zeta values; $54$ rational-ball product
containment instances; and $4$ overlapping independent-truncation checks.
Seven example enclosures are stored with exact rational centers and
radii, and outward-rounded decimal endpoints. These finite checks catch
implementation and normalization mistakes. They do not prove a
universal analytic theorem or constitute the object proof (L2).

For example, at $N=64$, $m=12$ and working precision $70$ bits,
the program returns the following outward-rounded decimal intervals.
The machine-readable file retains the exact rational ball radii.
\begin{center}
\small
\begin{tabular}{ll}
\toprule Quantity & Enclosure or exact value\\\midrule
$\zeta(2)$ & $[1.64493406684822,\ 1.64493406684823]$\\
$\zeta(3)$ & $[1.20205690315959,\ 1.20205690315960]$\\
$\zeta(1/2)$ & $[-1.46035450880959,\ -1.46035450880958]$\\
$\gamma$ & $[0.57721566490153,\ 0.57721566490154]$\\
$\zeta(-1)$ & $-1/12$ exactly\\
$\zeta(-2)$ & $0$ exactly\\
$\Re\zeta(1/2+14i)$ & $[0.02224114258984,\ 0.02224114262092]$\\
$\Im\zeta(1/2+14i)$ & $[-0.10325812328498,\ -0.10325812325391]$\\
\bottomrule
\end{tabular}
\end{center}
The complex example is a value enclosure, not a certificate for a zero.

From the extracted package root, run
\begin{verbatim}
python3 code/check_finite.py
python3 code/zeta_certificates.py --output verification/enclosures.json
python3 code/build_pdf.py
\end{verbatim}
The first two commands use only Python's standard library. The third
needs \texttt{pdflatex} and the ordinary LaTeX packages listed in the
preamble. Proposed Lean checking commands and the exact module roadmap
are in \texttt{lean/README.md}. The PDF build, exact checks, and visual
inspection are recorded separately from unperformed Lean verification.
